\documentclass[10pt,a4paper]{amsart}
\usepackage[margin=19mm]{geometry}
\usepackage{amsmath,amssymb,mathtools}
\usepackage[scr=rsfs,cal=euler]{mathalfa}
\usepackage{xfrac}
\usepackage[unicode,hidelinks,bookmarks=false]{hyperref}
\newcommand{\ie}{i.e.\ }
\newcommand{\cf}{cf.\ }
\newcommand{\resp}{respectively\ }
\newcommand{\Subsection}{\subsection}
\providecommand{\nobreakdash}{\nobreak\hbox{-}}
\setlength{\emergencystretch}{2em}
\multlinegap0pt
\usepackage{enumerate}
\usepackage{amssymb}
\newtheorem{lem}{Lemma}
\newcommand{\cH}{{\mathcal H}}
\newcommand{\cK}{{\mathcal K}}
\title{L\"{o}wner equations and de Branges--Rovnyak spaces}
\author{Michio Seto}
\date{Source: arXiv:2512.10368v4 (CC BY 4.0)}
\begin{document}
\maketitle

\noindent\emph{Source and licence.} L\"{o}wner equations and de Branges--Rovnyak spaces. Version arXiv:2512.10368v4 (CC BY 4.0), distributed by arXiv under the Creative Commons Attribution 4.0 International licence, http://creativecommons.org/licenses/by/4.0/, as recorded in the arXiv licence metadata for this exact version and shown on the abstract page https://arxiv.org/abs/2512.10368v4. Full licence notice: \texttt{licenses/arxiv-2512.10368-CC-BY-4.0.txt}.\par\smallskip

\noindent\textit{Editorial scope.} Counted unit: the article's lem lem:4-1 (source0.tex lines 244--260). The article's complete original proof of it is delivered with it (source0.tex lines 262--301, one \texttt{proof} environment of 40 lines). No mathematics is written, repaired or re-derived: the statement and the proof are the article's own lines, in the article's own notation, and no retained line is substituted or re-flowed.  No further source-local context is needed: every cross-reference of every retained line resolves inside the two delivered spans.  Nothing was omitted from any retained span, so the package carries no omission record.  No retained line contains a citation, so the wrapper carries no bibliography and no bibliography entry.  No retained line contains a figure environment, an image-inclusion command or drawing code, so no image is delivered and the package declares no asset dependency.  Editorial formatting only, in addition: an AMS \texttt{amsart} preamble that restates the article's own declarations and macros used by the retained lines, namely the environments \texttt{lem} on separate counters, together with the standard packages those lines need.  The retained environments print in this document as Lemma 1 in order of appearance, whereas the article prints the same items with the numbering of its own full document.  The complete native source of the article as distributed in the arXiv e-print for this exact version is bundled unchanged as \texttt{sources/190814/source0.tex}.
\begin{lem}\label{lem:4-1}
Let $\{f_n\}_n$ be a sequence in $\cK$. 
If 
\[
\int_X\|f_n(x,\cdot)-f_m(x,\cdot)\|_{\cH(x)}^2\ d\sigma(x)\to 0\quad(n,m\to \infty),
\] 
then, 
for almost all $x$,   
there exists a function $f(x, z)$ in $\cH(x)$ such that 
\[
\int_X\|f(x,\cdot)\|_{\cH(x)}^2\ d\sigma(x)<\infty
\]
and
\[
\int_X\|f_n(x,\cdot)-f(x,\cdot)\|_{\cH(x)}^2\ d\sigma(x)\to 0\quad(n\to \infty).
\] 
\end{lem}

\begin{proof}
Since this lemma is essentially the same as the Riesz--Fischer theorem, 
we give a sketch of the proof. 
It follows from the assumption that 
there exists a subsequence $\{f_{n_k}\}_k$ of $\{f_n\}_n$ such that 
\[
\int_X \|f_{n_{k+1}}(x,\cdot)-f_{n_k}(x,\cdot)\|_{\cH(x)}^2\ d\sigma(x)<\frac{1}{2^{2k}}. 
\]
Then, since 
\[
g(x)=\|f_{n_1}(x,\cdot)\|_{\cH(x)}+\sum_{k=1}^{\infty}\|f_{n_{k+1}}(x,\cdot)-f_{n_k}(x,\cdot)\|_{\cH(x)} 
\]
is square integrable by the monotone convergence theorem. 
In particular, $g(x)$ is finite for almost all $x$ in $X$. 
Hence,
we have that
\[
\|f_{n_{\ell}}(x,\cdot)-f_{n_k}(x,\cdot)\|_{\cH(x)}
\leq \sum_{j=k}^{\ell-1}\|f_{n_{j+1}}(x,\cdot)-f_{n_j}(x,\cdot)\|_{\cH(x)}\to 0\quad(k,\ell\to\infty)
\]
for almost all $x$ in $X$. 
Therefore, 
for almost all $x$,   
there exists a function $f(x, z)$ in $\cH(x)$ 
such that 
\[
\lim_{k\to \infty}\|f_{n_k}(x,\cdot)-f(x,\cdot)\|_{\cH(x)}=0. 
\]
Then, 
since $\|f_{n_k}(x,\cdot)\|_{\cH(x)}\leq g(x)$, 
by Lebesgue's dominated convergence theorem, we have 
\[
\int_X\|f(x,\cdot)\|_{\cH(x)}^2\ d\sigma(x)<\infty.
\]
Further, 
by Fatou's lemma, we have 
\[
\int_X\|f_n(x,\cdot)-f(x,\cdot)\|_{\cH(x)}^2\ d\sigma(x)\to 0\quad(n\to \infty).
\] 
\end{proof}

\end{document}
